% SmartNotes – preamble for kjeminotater (stil «Moderne», farger fra «Tavle»).
% Claude får IKKE lov til å endre denne; alt innhold skrives i dokumentkroppen.
\documentclass[11pt,a4paper]{article}

\usepackage[T1]{fontenc}
\usepackage[utf8]{inputenc}
\usepackage[norsk]{babel}
\usepackage[scaled=0.95]{helvet}
\renewcommand{\familydefault}{\sfdefault}
\usepackage{sfmath}
\usepackage{microtype}
\usepackage[a4paper,margin=2.3cm,headheight=14pt]{geometry}
\usepackage[parfill]{parskip}

% --- Matematikk ---
\usepackage{amsmath,amssymb,mathtools}
\usepackage{icomma}            % 9,81 skrives riktig i matematikkmodus
\usepackage{cancel}            % \cancel{} for forkorting
\usepackage{siunitx}
\sisetup{
  output-decimal-marker = {,},
  group-separator = {\,},
  group-minimum-digits = 5,
  exponent-product = \cdot,
  inter-unit-product = \cdot,
  per-mode = symbol,
  range-phrase = {\text{ til }},
  list-final-separator = {\text{ og }},
  list-pair-separator = {\text{ og }},
}
\usepackage[version=4]{mhchem} % formler og reaksjonslikninger: \ce{2 H2 + O2 -> 2 H2O}
\usepackage{chemfig}             % strukturformler: \chemfig{H-C(-[2]H)(-[6]H)-H}
\setchemfig{atom sep=2em,bond offset=1pt,double bond sep=2.5pt}

% Egne makroer (dokumentert i instruksene til Claude)
\DeclarePairedDelimiter{\abs}{\lvert}{\rvert}
\DeclarePairedDelimiter{\norm}{\lVert}{\rVert}
\newcommand{\dd}{\mathop{}\!\mathrm{d}}
\newcommand{\dv}[2]{\frac{\mathrm{d}#1}{\mathrm{d}#2}}
\newcommand{\ddv}[2]{\frac{\mathrm{d}^2#1}{\mathrm{d}#2^2}}
\newcommand{\pdv}[2]{\frac{\partial#1}{\partial#2}}
\newcommand{\vb}[1]{\mathbf{#1}}
\newcommand{\vu}[1]{\hat{\mathbf{#1}}}
\newcommand{\enhet}[1]{\,[\mathrm{#1}]}

% --- Figurer og tegninger ---
\usepackage{graphicx}
\usepackage{float}
\usepackage[labelfont=bf,font=small]{caption}
\usepackage{tikz}
\usetikzlibrary{arrows.meta,calc,angles,quotes,patterns,positioning,
  decorations.pathmorphing,decorations.markings,shapes.geometric,intersections,babel}
\usepackage{pgfplots}
\pgfplotsset{compat=1.18}
\usepackage[european,siunitx]{circuitikz}

% --- Tabeller og lister ---
\usepackage{booktabs,array,tabularx}
\usepackage{enumitem}

% --- Bokser ---
\usepackage{xcolor}
\usepackage[most]{tcolorbox}
% Tavle: tavlegrønn som blekkfarge og krittgul aksent (mørknet litt for utskrift)
\definecolor{snblue}{HTML}{2E4A3F}
\definecolor{sngreen}{HTML}{2E4A3F}
\definecolor{snorange}{HTML}{C2620A}
\definecolor{sngray}{HTML}{59625D}
\definecolor{snyellow}{HTML}{A8821C}
\tcbset{snbase/.style={enhanced,breakable,boxrule=0.5pt,arc=3pt,
  left=8pt,right=8pt,top=8pt,bottom=6pt,fonttitle=\bfseries\footnotesize,
  attach boxed title to top left={xshift=8pt,yshift=-2.2mm},
  boxed title style={boxrule=0pt,arc=2pt,left=4pt,right=4pt,top=1pt,bottom=1pt},
  before skip=12pt,after skip=10pt}}

\newtcolorbox{definisjon}[1][]{snbase,colback=white,colframe=snblue!45,colbacktitle=snblue,coltitle=white,
  title={Definisjon\ifx\relax#1\relax\else: #1\fi}}
\newtcolorbox{formel}[1][]{snbase,colback=white,colframe=sngreen!45,colbacktitle=sngreen,coltitle=white,
  title={\ifx\relax#1\relax Viktig\else #1\fi}}
\newtcolorbox{eksempel}[1][]{snbase,colback=white,colframe=sngray!45,colbacktitle=sngray,coltitle=white,
  title={Eksempel\ifx\relax#1\relax\else: #1\fi}}
\newtcolorbox{oppgave}[1][]{snbase,colback=white,colframe=snblue!45,colbacktitle=snblue,coltitle=white,
  title={Oppgave\ifx\relax#1\relax\else\ #1\fi}}
\newtcolorbox{losning}{snbase,colback=snblue!3,colframe=snblue!25,colbacktitle=snblue!15,coltitle=snblue!90!black,title={Løsning}}
\newtcolorbox{husk}{snbase,colback=white,colframe=snyellow!45,colbacktitle=snyellow,coltitle=white,
  title={Husk}}
\newtcolorbox{merknad}{snbase,colback=white,colframe=snorange!45,colbacktitle=snorange,coltitle=white,
  title={Merknad fra SmartNotes}}

\newcommand{\uleselig}[1]{\textcolor{sngray}{\textit{#1}}\textsuperscript{\textcolor{snorange}{\,?}}}
\newcommand{\sidenotat}[1]{\marginpar{\footnotesize\raggedright #1}}

% Figurer klippet ut fra originalbildet. \notedir settes av SmartNotes.
\newcommand{\notedir}{.}
\newcommand{\originalfigur}[3][0.6]{%
  \begin{figure}[H]\centering
  \includegraphics[width=#1\linewidth,height=0.45\textheight,keepaspectratio]{\notedir/#2}%
  \ifx\relax#3\relax\else\caption{#3}\fi
  \end{figure}}
\newcommand{\manglerfigur}[1]{%
  \begin{center}\fbox{\parbox{0.6\linewidth}{\centering\small\textcolor{sngray}{[Figur mangler: #1]}}}\end{center}}

% --- Sidehode ---
\usepackage{fancyhdr}
\pagestyle{fancy}
\fancyhf{}
\renewcommand{\headrulewidth}{0.4pt}
\newcommand{\snfag}{}
\newcommand{\snkapittel}{}
\fancyhead[L]{\small\textcolor{sngray}{\snfag}}
\fancyhead[R]{\small\textcolor{sngray}{\snkapittel}}
\fancyfoot[C]{\small\thepage}

% Ingen nummerering av overskrifter i notater
\setcounter{secnumdepth}{0}
\setcounter{tocdepth}{1}

% Tittelblokk for ett notat
\newcommand{\notehode}[3]{%
  \begingroup
  \noindent{\LARGE\bfseries #1\par}\vspace{2pt}
  \noindent{\small\textcolor{sngray}{#2\ifx\relax#3\relax\else\ \textbullet\ #3\fi}\par}
  \vspace{4pt}\noindent\textcolor{snblue}{\rule{\linewidth}{1pt}}\par
  \endgroup\vspace{6pt}}

\usepackage{titlesec}
\titleformat{\section}{\Large\bfseries}{}{0pt}{}[\vspace{-6pt}\textcolor{sngray!40}{\rule{\linewidth}{0.4pt}}]
\titleformat{\subsection}{\large\bfseries\color{snblue}}{}{0pt}{}
\renewcommand{\notehode}[3]{%
  \begingroup
  \noindent{\small\bfseries\textcolor{snblue}{#2}\par}\vspace{2pt}
  \noindent{\huge\bfseries #1\par}\vspace{3pt}
  \noindent{\small\textcolor{sngray}{#3}\par}
  \endgroup\vspace{10pt}}

\usepackage[hidelinks]{hyperref}
\usepackage{bookmark}
